% !TEX root = ../../../main.tex
% Appendix entry: detection.summation_threshold -- "Summation threshold" (Appendix A).
\entryfield{Algorithm}
Stage~3 of the \texttt{dehshibi\_spikes} template, after \blockref{preprocessing.wavelet_transform}
and \blockref{detection.wavelet_summation}. \code{_run} in
\code{Adapters/detection_summation_threshold.py} calls \code{detect_from_omega} in
\code{Working/Detection/analysis/dehshibi_authors.py}, which ports the authors' MATLAB
(\code{iApplyWavelet.m} after $\Omega$, \code{iClusterWaveletROI.m}, \code{iApplyEnvelope.m},
\code{iLocateSpike.m}) rather than the printed Algorithms~1--4. Let $p =
\mathrm{round}(\texttt{min\_separation\_s}\cdot f_s)$ (or \emph{n\_p} if non-zero) and
$\varepsilon = \texttt{epsilon\_factor}$. The score is cut into the same chunks as stage~1
(\code{chunks_for}: fixed windows of \emph{window\_s}, or \code{split_signal} with
\emph{slice\_by\_state}). A chunk shorter than 4 samples, or whose score is not finite
throughout, is skipped and counted. In each remaining chunk, with score $\Omega$ and signal $x$:
\begin{enumerate}
\item[B] \emph{Candidate regions} (\code{wavelet_regions}). Peaks of $\Omega$ and peaks of
  $\max\Omega - \Omega$ (valleys) by a port of MATLAB's \code{findpeaks}: prominence at least
  $\varepsilon(\max\Omega - \min\Omega)$, half-prominence width at least $p/2$, then, tallest
  first, no two within $p$ samples. Peaks and valleys are sorted together and paired
  consecutively, (1st, 2nd), (3rd, 4th), \dots; an odd count gets the chunk's last sample as
  its final mark, and that last pair is shortened to at most the mean length of the others.
  Every region's start is then moved earlier by $\lfloor s/2\rfloor$, where $s$ is the
  floored standard deviation of the region lengths, or by $\lfloor p/2\rfloor$ if $s > p$.
\item[C/D] \emph{Spike or pseudo-spike} (\code{cluster_regions}). A region shorter than $p$
  samples is dropped. In the rest, take the interior local minima and maxima of $x$ (scipy
  \code{find_peaks}); the region is a spike candidate (C) if the lowest minimum lies below
  both end values or the highest maximum above both; otherwise it is a pseudo-spike (D).
  D is reported and never used again, as in the authors' code.
\item[R] \emph{Envelope regions} (\code{envelope_regions}), from the signal alone. $L =
  \mathtt{del2}(x)$, a quarter of the second difference with the ends extrapolated; upper and
  lower envelopes are not-a-knot splines through the local maxima and minima of $L$ at least
  $p$ apart (MATLAB's \code{envelope(L, p, 'peak')}); $m$ is their mean. The minima of $m$ and
  its maxima less the first are sorted together and paired consecutively; each pair's strength
  is $\delta = |m(\mathrm{end}) - m(\mathrm{start})|$.
\item[S] \emph{Locate} (\code{locate_spikes}). Keep the envelope regions with $\delta \ge
  \bar\delta - \sigma_\delta$ if $\sigma_\delta < \bar\delta$, else $\delta \ge \bar\delta$
  (R\_kept). Widen each C region to the hull of the kept envelope regions it overlaps (searched
  between the envelope starts nearest its two ends), then merge consecutive results that
  overlap. If no C region overlaps any envelope region, the C regions are returned unchanged.
\end{enumerate}
Region indices are inclusive in the core and converted to half-open spans at the adapter.

\emph{Where the port follows the code and not the paper} (from the clause table in
\code{docs/prompts/fixup/reports/J-dehshibi-vs-the-paper.md}, \S4): extrema need minimum
distance and width as well as prominence; region starts are moved back; regions need at least
$p$ samples (the paper says more than 30); extrema are local, not spline-interpolated; the
envelope is a peak spline of $L$ itself, with no analytic signal; Algorithm~3's pairing is
consecutive, so no region can end before it starts; Algorithm~4 uses C only, merges instead
of emitting nested spans, and has no length filter.

\emph{How it connects.} The block needs a per-sample \texttt{Scores} and the signal $x$, which
every block receives; the envelope branch never looks at the score. So any per-sample score
can stand in for $\Omega$ (\code{test_any_per_sample_score_can_stand_in_for_omega}).

\entryfield{Assumptions}
\begin{itemize}
\item \emph{Units.} Not unit-aware, and the spans do not depend on the unit: $\varepsilon$ is
  relative to the score's range, the C/D test compares the signal with itself, and the
  envelope cut is relative to the envelope regions' own strengths. Only the span scores (the
  signal's range) are in the recording's stored unit.
\item \emph{Sampling rate.} Lengths are in seconds and converted with $f_s$. The authors ran at
  1\,Hz. At 10\,Hz the envelope stage, which works on sample-to-sample curvature, sees a
  different signal, so spans differ though the same events are found
  (\code{test_lengths_are_seconds_not_samples}).
\item \emph{Windows.} \emph{window\_s} and \emph{slice\_by\_state} must equal those given to the
  transform; nothing checks it.
\item \emph{Missing values.} A window whose score holds a NaN is skipped (counted in
  \code{n_chunks_skipped}). The block does not check $x$ itself: a NaN in $x$ under a finite
  score from some other block would pass into \code{find_peaks} and the splines with
  \unverified{unpredictable results}. Behind the transform this cannot happen, since a NaN in
  $x$ makes that window's score NaN.
\end{itemize}

\entryfield{Outputs}
A \texttt{SpanSet}, one half-open span per spike, label \texttt{spike}, score the range
($\max - \min$) of the signal inside it. \code{meta} holds \code{pseudo_spikes} (D, half-open),
the counts \code{n_spikes}, \code{n_pseudo_spikes}, \code{n_chunks}, \code{n_chunks_skipped},
\code{n_candidates_B}, \code{n_wavelet_C}, \code{n_envelope_R}, and \code{funnel}: for each
stage B, C, D, R, R\_kept, S its label, count and regions, plus the chunks and the indices of
the peaks and valleys of $\Omega$. The Analyse page draws the funnel
(\code{webui/client/src/analyse/DetectorFunnel.tsx}).

\entryfield{Cost}
Per window: the peak search, prominence and width are about linear in the window; the
distance filter is quadratic in the number of peaks; the splines are linear. Overall roughly
linear in the span for a fixed window. No estimate or span cap is declared; the transform
upstream caps the template at 320{,}000 samples.

\entryfield{Failure modes}
\begin{itemize}
\item \emph{Unselective on real traces.} The prominence is relative to each window's own
  range, so every window produces extrema; and the C test passes whenever the signal inside a
  region dips below or rises above both its ends, which a noisy trace nearly always does. On
  M2\_aug CH0 all 38 candidate regions in four windows passed it (fixup-J report, \S6).
\item \emph{A candidate can vanish in step S.} When at least one C region meets an envelope
  region, any C region that meets none is dropped. Only when none meets any are they all kept.
\item \emph{Merging.} Step S's merge sets the earlier region's end to the later one's without
  taking the larger; were a later region ever to end inside an earlier one, the merged span
  would be cut short. \unverified{Whether this occurs in practice} was not checked.
\item \emph{Window mismatch} with the transform (above).
\item \emph{The tail.} The last window can be up to 1.5 windows long, so its relative thresholds
  cover more signal.
\end{itemize}

\entryfield{Validation}
Against the authors' own code: \code{tests/test_dehshibi_authors.py} and
\code{tests/test_adapter_summation_threshold.py} give the port the $\Omega$ MATLAB computed on
three synthetic recordings (\code{tests/fixtures/dehshibi/reference.json}) and require every
region set (B, C, D, R) and the final spikes to match MATLAB exactly. So the port is faithful.

Against known events (\code{scripts/fixup_j_figure.py}, results in
\code{webui/screenshots/fixup/J/known-events.json}; an event is found when one span covers at
least half of it and is no more than three times as long):
\begin{itemize}
\item a 12{,}000\,s synthetic recording with 20 injected events of 0.5--6\,mV
  (\code{scripts/fixup_j_synthetic.py}): 17 found, 19 spans, none on no event. The three missed
  had amplitudes 2.7, 2.8 and 4.5\,mV, so the misses are not the small events;
\item the three MATLAB reference recordings: 5 of 6, 3 of 4 and 8 of 8 by that criterion
  (the tests, counting any overlap, require at least 5, 4 and 8);
\item M2\_aug CH0, 683 ten-minute windows a researcher labelled, each run in the 3000\,s window
  it falls in: a span in 130 of 143 \emph{interesting} windows (91\,\%) and in 466 of 540
  \emph{not interesting} ones (86\,\%); 2{,}846 spans covering 53\,\% of the time. Before the
  port the figures were 13\,\% and 12\,\%, with 15\,\% of the time covered.
\end{itemize}
The synthetic says the port detects the paper's kind of event on a quiet baseline. The labels
say that on this recording it cannot tell an interesting window from an uninteresting one. The
paper reports 76\,\% true positives and 16\,\% false positives on one 36{,}000\,s piece of its
own data \citep{dehshibi2021electrical}.

\entryfield{References}
\citet{dehshibi2021electrical}; the authors' code, Zenodo doi:10.5281/zenodo.3997031; the
decisions behind the port are \code{docs/prompts/fixup/QUESTIONS.md}, Q28--Q33.

\entryfield{Tests}
\code{tests/test_adapter_summation_threshold.py} (types, seconds, old parameter names, spikes
and pseudo-spikes against MATLAB, labels and scores, no nested spans, the funnel, a stand-in
score, the sampling rate, refusals, all-NaN scores give no spans);
\code{tests/test_dehshibi_authors.py} (each stage of the core, two windows, state slicing);
\code{tests/test_template_dehshibi.py} (the template end to end, written to the database).
